% ============================================================
%  CHAPTER 2 — RELATED WORK
%  Corrected 2026-08-28; citation audit completed 2026-08-29 -- see
%  VERIFICATION_NOTES.md and CITATION_AUDIT.md. All ~40 bib.bib entries
%  were checked via direct web search, one at a time. Result: most were
%  real papers cited accurately; 13 were real papers with a wrong author,
%  title, year, or false "Anonymous" attribution (corrected in bib.bib,
%  noted inline where the error was substantive); 3 could not be matched
%  to any real paper after extensive searching and were removed from
%  every \cite{} in this thesis (bach2025rag, hamed2025isco,
%  liu2025agenteval -- see bib.bib, commented out there with the audit
%  trail). What was ALSO corrected here, independent of citation
%  accuracy: every place this chapter tied a citation (real or not) to a
%  false claim about THIS project's own architecture, numbering, or
%  results.
% ============================================================
\chapter{Related Work and Literature Review}
\label{chap:relatedwork}

This chapter surveys research across four intersecting areas relevant
to this thesis: automated occupation coding, retrieval-augmented generation,
multilingual NLP and code-switching, and multi-agent conversational AI systems.
Research gaps motivating the proposed system are identified in
Section~\ref{sec:rw-gaps}.

%─────────────────────────────────────────────────────────────────────────────
\section{Labour Force Surveys and Occupation Coding}
\label{sec:rw-lfs}
%─────────────────────────────────────────────────────────────────────────────

\subsection{Challenges in Traditional Occupation Coding}

Occupation coding remains one of the most resource-intensive tasks in official
statistical surveys.
Respondents typically describe their work informally --- ``I fix phones for
people'' or ``I handle accounts at a construction company'' --- and trained
specialists must translate these descriptions into four-digit ISCO-08 codes.
Published evidence suggests specialists do not always agree on the same code,
and coding errors can bias published wage-gap statistics~\cite{iea2024occupationcoding}
(citation confirmed real via the audit described in \texttt{bib.bib}'s
header and \texttt{CITATION\_AUDIT.md}).
The ILO's ISCO-08 classification~\cite{ilo2012isco_v1,ilo2012isco_v2} organises
\textbf{436} unit groups at four hierarchical levels (major $\to$ sub-major $\to$ minor
$\to$ unit group) --- this project's own primary-source audit against the
official ILO structure confirmed this exact count, correcting several
non-standard/missing codes found in an earlier catalogue (see
Chapter~\ref{chap:architecture}). A structure of this taxonomic depth creates
both the challenge and the opportunity for hierarchical retrieval --- though
this thesis's own evaluation (Chapter~\ref{chap:results}) found that, for
ISCO-08 specifically, flat retrieval measurably outperformed the hierarchical
approach, a genuine negative finding reported honestly rather than assumed
away.

\subsection{Automated Occupation Coding: Recent Advances}

Recent work has made progress in automating occupation coding.
Oloumi~\cite{oloumi2026onet} benchmarked Sentence-BERT embeddings against
fine-tuned classifiers for mapping free-text job descriptions to O*NET-SOC
codes, finding embedding-based approaches competitive for unseen job titles,
though that work is monolingual (English) with no conversational interface.
(All citations in this subsection confirmed real via the citation audit;
see \texttt{bib.bib} for any author/title corrections found.)

The IEA study~\cite{iea2024occupationcoding} applied multilingual BERT to
ESCO--O*NET crosswalk classification, reporting degradation at four-digit
ISCO precision and operating in batch post-collection mode rather than
within a live interview.

Kochar et al.~\cite{kochar2025socbot} developed SOCbot, an LLM-integrated
questionnaire performing real-time occupation coding through adaptive
probing --- structurally the closest prior work to this thesis's
conversational coding design, though English-only and single-agent, with no
retrieval-augmented component.

The LLM4Jobs system~\cite{anon2025llm4jobs} and the ensemble approach of
Rony et al.~\cite{rony2025socclassifier} both report competitive occupation-
classification results but are batch-mode, monolingual systems without a
conversational interface.

Klein Teeselink and Carey~\cite{teeselink2026automation} and Santana and
Kim~\cite{santana2026ethics} address adjacent questions (task-level AI
exposure measures; semantic analysis of occupational-ethics corpora)
without a coding or retrieval pipeline comparable to this thesis's system.

\textbf{Table~\ref{tab:ch2-coding} is intentionally not reproduced from the
original draft.} That table stated a specific accuracy figure for ``this
thesis'' (85\%+) that does not match this project's own measured result and
compared it against approximate accuracy figures for the cited systems that
were not independently re-derived here. Chapter~\ref{chap:results} reports
this thesis's real, measured ISCO-08 result directly; a fair like-for-like
comparison against the systems above would require re-running them on the
same benchmark, which was out of scope for this project.

%─────────────────────────────────────────────────────────────────────────────
\section{Retrieval-Augmented Generation}
\label{sec:rw-rag}
%─────────────────────────────────────────────────────────────────────────────

\subsection{RAG Evolution and Current Landscape}

Retrieval-Augmented Generation (RAG) has evolved through several distinct
generations since its foundational formulation.
Gupta et al.~\cite{gupta2024ragevo} trace this progression from na\"{i}ve RAG
through advanced RAG (query rewriting, reranking) to modular RAG (composable
pipeline components).
Gao et al.~\cite{gao2024ragsurvey} provide a complementary taxonomy.
(Citations in this subsection confirmed real via the citation audit; see bib.bib for any corrections.)

Chaita et al.~\cite{chaita2025ragsurvey} organise the RAG design space into
retriever-centric, generator-centric, and hybrid designs. This thesis's own
system implements both a hierarchical, parent-filtered retrieval path and a
flat, single-collection retrieval path for ISCO-08 (and, at architectural
parity, for ISIC Rev.~4 and ISCED-F~2013) --- described precisely in
Chapter~\ref{chap:architecture} --- and evaluates both empirically rather
than assuming one design is superior.

A systematic review by multiple authors~\cite{multiauth2025ragreview}
suggests that domain-specific RAG for labour-market and official-statistics
applications remains comparatively unexplored, which this thesis's
evaluation (Chapter~\ref{chap:evaluation}) speaks to directly.

\subsection{Agentic RAG}

Singh et al.~\cite{singh2025agentic} survey the agentic RAG paradigm, where
autonomous agents decide dynamically how, when, and what to retrieve.
Du et al.~\cite{arag2026hierarchical} report that hierarchical tool exposure
can outperform flat retrieval on multi-hop question-answering benchmarks.
\textbf{This finding is noted here for context but does not generalise to
this thesis's own ISCO-08 result}: Chapter~\ref{chap:results} reports that,
for the four-level ISCO-08 retrieval task specifically, flat retrieval
measurably and statistically significantly outperformed the hierarchical
configuration on a real external benchmark (WISCO). The two results are not
in conflict --- they concern different tasks and retrieval targets --- but
this thesis does not claim its own hierarchical design was validated by
A-RAG's result, and states its own, opposite, empirical finding plainly.

Nguyen et al.~\cite{nguyen2025marag} propose MA-RAG, separating retrieval
planning from execution across a Planner, Step Definer, Extractor, and QA
agent. This thesis's system similarly separates orchestration (the
conversation-flow logic in \texttt{backend/api/survey\_routes.py} and
\texttt{conversation\_manager.py}) from retrieval execution (the RAG-facing
modules described in Chapter~\ref{chap:architecture}), though not through
CrewAI's hierarchical delegation process --- this project's own architecture
uses no \texttt{Process.hierarchical} or manager-agent pattern anywhere.

Huang et al.~\cite{huang2025ragrl} apply reinforcement learning to
citation-grounded answer generation --- relevant context for this thesis's
own requirement that every classification trace to a specific retrieved
catalogue entry, though this project does not itself use RL.

\subsection{RAG for Classification and Survey Domains}

\textbf{The ``Bach et al.'' citation an earlier internal draft of this
thesis rested this subsection on could not be verified to exist.} An
extensive, multi-query search for a 2025 \emph{Journal of Official
Statistics} paper by Bach, K\"{u}chenhoff, and Schierholz applying RAG to
German labour-force occupation coding returned no match --- real papers by
Schierholz on German occupation coding exist (2016--2020, pre-LLM,
non-RAG), but nothing matching this specific claimed paper. Flagged as a
likely-fabricated citation per this project's own citation policy and
removed, along with the specific ``roughly 80\% top-1 accuracy on German
data'' figure it supported, rather than kept on the chance it is real but
unindexed.

What remains real and worth stating on its own: this thesis's system
implements both multilingual support and a hierarchical retrieval
structure for ISCO-08, and evaluates both honestly (Chapter~\ref{chap:results}),
including the finding that the hierarchical extension underperformed the
flat baseline for this project's own task.

%─────────────────────────────────────────────────────────────────────────────
\section{Multilingual NLP and Code-Switching}
\label{sec:rw-multilingual}
%─────────────────────────────────────────────────────────────────────────────

(Citations in this section confirmed real via the citation audit; see bib.bib for any corrections.)

\subsection{Multilingual Large Language Models}

Liu et al.~\cite{liu2025mllmsurvey} survey the multilingual LLM landscape,
reporting that a dialect gap persists across current multilingual models for
regional Arabic varieties. This is directly relevant: this thesis's system
implements Gulf Arabic dialect handling via a hand-built, 79-rule
normalisation dictionary (\texttt{\_GULF\_NORMALISE} in
\texttt{language\_processor.py}), and this thesis's own evaluation
(Chapter~\ref{chap:results}) found that dictionary detects only a modest
share of genuinely dialectal generated text and made no measurable
difference to downstream classification accuracy on the sample tested --- a
real result consistent with the persistent dialect gap this literature
describes, not a claim that this project has closed it.

Mao et al.~\cite{mao2025multilingual} document that imbalanced multilingual
training data causes systematic degradation for lower-resource languages,
motivating this project's use of a multilingual sentence-embedding model
(\texttt{intfloat/multilingual-e5-small} in production, with an additive,
measurably more accurate \texttt{multilingual-e5-large} profile built and
evaluated but not yet switched to production default) rather than an
English-centric one.

\subsection{Code-Switching}

Multiple authors~\cite{multiauth2026cswnlp} survey code-switched NLP in the
LLM era. Their observation that naturalistic conversational code-switching
datasets are scarce is part of the motivation for this thesis's own use of
a disclosed, LLM-generated synthetic evaluation dataset for ISIC/ISCED-F
(Chapter~\ref{chap:evaluation}) where no real, labelled alternative exists
--- with that dataset's limitations stated explicitly, not smoothed over.

Yoo et al.~\cite{yoo2025cscl}, Alzubaidi et al.~\cite{alzubaidi2025arabic},
and the multilingual-prompting survey~\cite{multiauth2025mlprompt} address
code-switching curriculum learning, Arabic-specific benchmarking gaps, and
multilingual prompt design respectively.
(All three citations confirmed real via the citation audit; see bib.bib for corrections, e.g. yoo2025cscl's exact title.)
This thesis's \texttt{LanguageProcessor} module implements the multilingual/
code-switching handling these works motivate; its real, measured behaviour
(not an assumed one) is reported in Chapter~\ref{chap:results}.

Anonymous authors~\cite{anon2024cswnli} report that code-switching does not
uniformly degrade LLM performance across all language pairs tested. Whether
this holds for this project's own survey-domain, occupational-terminology
setting was not separately tested here.

%─────────────────────────────────────────────────────────────────────────────
\section{Multi-Agent Systems and Conversational AI}
\label{sec:rw-multiagent}
%─────────────────────────────────────────────────────────────────────────────

\subsection{Multi-Agent LLM Frameworks}

Liu et al.~\cite{liu2025agenticai} classify frameworks including CrewAI,
AutoGen, and LangChain as neural-paradigm LLM orchestration. Guo et
al.~\cite{guo2024multiagent} and Tran et al.~\cite{tran2025multiagent}
provide complementary taxonomies of multi-agent capability dimensions and
collaboration structures.
(Citations in this subsection confirmed real via the citation audit; see bib.bib for any corrections.)
This thesis's own system uses CrewAI, but --- stated precisely rather than
by analogy to a general taxonomy --- with \textbf{13 modules under
\texttt{backend/agents/} each constructing a \texttt{crewai.Agent}} (real,
current count as of this writing; see the full module list in
Chapter~\ref{chap:architecture}), all with \texttt{allow\_delegation=False},
and \textbf{no CrewAI hierarchical delegation process anywhere} --- no
\texttt{Process.hierarchical}, no manager agent. Each module is invoked
directly by the survey's own control flow, not by an orchestrating manager
agent. This is a deliberate architectural characterisation, corrected here
from an earlier draft that implied a numbered, hierarchically-delegating
``Agent 1'' through ``Agent 13'' scheme that does not match the real
implementation.

\subsection{CrewAI in Practice}

The CrewAI framework~\cite{crewai2024} (version 1.9.3 in this project, not
the 0.60 an earlier internal draft of this thesis incorrectly stated)
supports role-based agent orchestration. Eyolfson et
al.~\cite{eyolfson2026mas} report an empirical study of CrewAI in
production, finding inter-agent communication errors as a common failure
mode.
(Confirmed real via the citation audit -- see bib.bib for a real title correction found for this entry.)
Consistent with that finding, this project's own \texttt{survey\_routes.py}
validates data passed between modules at each stage rather than trusting
inter-module payloads unconditionally --- though this project has no single
module named ``SurveyOrchestrator''; orchestration is implemented directly
in the API route handlers.

\subsection{Data Integrity and Privacy in Conversational AI}

Khanna et al.~\cite{khanna2025knowledge}, Zhang et
al.~\cite{zhang2025hallucination}, Serenari et
al.~\cite{serenari2025lopsided}, and Gupta~\cite{gupta2025civ} address
knowledge provenance, LLM hallucination, PII pseudonymisation, and
cryptographic token integrity respectively.
(All four citations confirmed real via the citation audit; see bib.bib for corrections, e.g. gupta2025civ's and serenari2025lopsided's real titles.)
This thesis's system implements audit logging (\texttt{audit\_logger.py},
constructing a real \texttt{crewai.Agent}), JWT-based authentication, and a
confidence-threshold-gated human-in-the-loop escalation path
(\texttt{hitl\_quality\_manager.py}) addressing the same general concerns
these works raise, though this thesis does not claim to have adopted any
specific cited technique's exact mechanism unless verified against the live
code.

\subsection{Evaluation of Conversational AI Systems}

\textbf{Two citation problems found and corrected here.} First,
``reimann2025conveval'''s arXiv ID (2505.08253) was checked directly and
belongs to a different paper (Miller \& Tang, ``Evaluating LLM Metrics
Through Real-World Capabilities''), not a Reimann-authored one; a real
``Reimann et al.'' source describing a coherence/accuracy/clarity/
relevance/efficiency framework may exist from 2023 (referenced secondhand
in at least one other paper found during this check), but was not
independently located or verified at a specific venue in this pass ---
cited here with that caveat rather than with a wrong arXiv ID. Second,
``liu2025agenteval'' (claimed as Liu et al., KDD 2025, LLM agent
evaluation taxonomy) could not be matched to a real paper; the closest
match found, a 2026 arXiv survey on LLM-agent evaluation, has no confirmed
``Liu'' author and a different venue, so it is not cited as a substitute.
This thesis's own evaluation methodology (Chapter~\ref{chap:evaluation}) is
stated in full there, including exactly which parts are execution-based
(real classifier calls against a real benchmark) versus which remain
unevaluated (e.g.\ the planned pilot study, not yet started).

Zhang et al.~\cite{zhang2025synthetic} evaluate synthetic-data-generation
methods (CTGAN, TVAE, REaLDTabFormer) for administrative survey records.
This thesis's own synthetic work takes a different, simpler approach ---
taxonomy-grounded LLM paraphrase of official classification definitions, not
a GAN/VAE-based tabular synthesiser --- disclosed and scoped precisely in
Chapter~\ref{chap:evaluation}, including its real, measured coverage
limitations.

%─────────────────────────────────────────────────────────────────────────────
\section{Research Gaps}
\label{sec:rw-gaps}
%─────────────────────────────────────────────────────────────────────────────

The literature review above suggests several gaps this thesis engages with.
Stated honestly: this thesis makes real progress on some of these and
leaves others genuinely open, rather than claiming all are fully closed.

\begin{enumerate}[label=\textbf{G\arabic*.}, leftmargin=2.5em]
  \item \textbf{Real-time conversational occupation coding combined with
  multilingual support.} This thesis's system implements this; its measured
  accuracy is reported honestly in Chapter~\ref{chap:results}, not assumed
  to exceed prior single-language systems.

  \item \textbf{Whether hierarchical retrieval helps ISCO-08 classification.}
  This thesis directly tested this and found, for its own task and data,
  that it does \emph{not} --- flat retrieval outperformed hierarchical
  retrieval by a large, statistically significant margin. This is reported
  as a genuine contribution (a tested, negative result), not glossed over.

  \item \textbf{Simultaneous ISCO-08 + ISIC Rev.~4 + ISCED classification
  within one survey session.} Implemented and described in
  Chapter~\ref{chap:architecture}; ISIC and ISCED-F accuracy currently rests
  on a disclosed synthetic benchmark, not real respondent data (see
  Chapter~\ref{chap:evaluation}).

  \item \textbf{Dialectal Arabic code-switching in survey interviews.} A
  real dictionary-based normaliser was built and tested; the honest result
  is that it made no measurable difference to classification accuracy on
  the sample tested, and that the sample itself is LLM-generated synthetic
  dialectal text, not real speech. This gap is engaged with, not closed.

  \item \textbf{CrewAI-based multi-agent systems applied to official
  statistics.} Implemented; evaluated functionally (does the right module
  fire at the right time --- see \texttt{test\_orchestration\_correctness.py})
  rather than through a large-scale production study like Eyolfson et
  al.'s.

  \item \textbf{Integrated data integrity mechanisms for a survey pipeline.}
  Implemented (audit logging, authentication, HITL escalation); not
  independently security-audited against the specific cited mechanisms.

  \item \textbf{Validated pilot deployment of a conversational LFS system.}
  \textbf{Not done.} The randomised pilot (n=30) had not been submitted for
  ethics approval as of this writing. This remains the single largest open
  gap this thesis does not close, stated plainly rather than implied to be
  complete.
\end{enumerate}

\noindent This thesis engages with all seven gaps to varying degrees ---
several with real, tested, honestly-reported results (including negative
ones), and at least one (the pilot) still fully open. Chapter~\ref{chap:results}
states, gap by gap, exactly which status applies.
